% =====================================================================
%  nebula_preamble.tex  ·  Nebula Courses · plantilla fija v3
%  Portada Nebula + sistema fijo de cajas semanticas + footer obligatorio.
%  El documento que incluye este fichero define antes:
%     \def\nebulalang{es}   o   \def\nebulalang{en}
% =====================================================================

\usepackage[utf8]{inputenc}
\usepackage[T1]{fontenc}
\usepackage{lmodern}

% RULE-013: la carga condicional de paquetes se hace con \ifx de TeX puro.
\providecommand{\nebulalang}{es}
\def\nb@es{es}
\ifx\nebulalang\nb@es
  \usepackage[spanish,es-noquoting,es-noshorthands]{babel}
\else
  \usepackage[english]{babel}
\fi

\usepackage{ifthen}
\newcommand{\nbIfES}[2]{\ifthenelse{\equal{\nebulalang}{es}}{#1}{#2}}

\usepackage{geometry}
\usepackage{microtype}
\usepackage{amsmath,amssymb,amsfonts,mathtools}
\usepackage{cancel}
\usepackage{siunitx}
\ifx\nebulalang\nb@es
  \sisetup{group-separator={\,},group-minimum-digits=4,output-decimal-marker={,}}
\else
  \sisetup{group-separator={\,},group-minimum-digits=4,output-decimal-marker={.}}
\fi
\usepackage{array,booktabs,tabularx,longtable}
\usepackage{enumitem}
\usepackage{multicol}
\usepackage{graphicx}
\usepackage[export]{adjustbox}
\usepackage{tikz}
\usetikzlibrary{calc,arrows.meta,positioning,decorations.pathmorphing,patterns,
  angles,quotes,shapes.geometric,intersections,through,3d,decorations.markings}
\usepackage{pgfplots}
\pgfplotsset{compat=1.16}
\usepackage{xcolor}
\usepackage[most]{tcolorbox}
\usepackage{fancyhdr}
\usepackage{titlesec}
\usepackage{needspace}
\usepackage{etoolbox}
\usepackage{float}
\usepackage[hidelinks]{hyperref}

\geometry{top=2.4cm,bottom=2.8cm,left=2.2cm,right=2.2cm}
\setlength{\parindent}{0pt}
\setlength{\parskip}{6pt}
\setlist[itemize]{topsep=3pt,itemsep=2pt,leftmargin=1.2em}
\setlist[enumerate]{topsep=3pt,itemsep=2pt,leftmargin=1.4em}
\renewcommand{\arraystretch}{1.2}

% =====================================================================
%  REGLAS DE SEGURIDAD DE MAQUETACION
% =====================================================================
% RULE-001/022: ninguna tabla excede el area de texto.
\newcommand{\nbtable}[1]{%
  \par\medskip
  \noindent\resizebox{\ifdim\width>\linewidth \linewidth\else \width\fi}{!}{#1}%
  \par\medskip}
% RULE-002: una ecuacion larga se encoge en vez de salirse del margen.
\newcommand{\nbeq}[1]{%
  \begin{equation*}\resizebox{\ifdim\width>\linewidth \linewidth\else \width\fi}{!}{$\displaystyle #1$}\end{equation*}}
% RULE-004: ningun encabezado huerfano al pie de pagina.
\pretocmd{\section}{\Needspace{7\baselineskip}}{}{}
\pretocmd{\subsection}{\Needspace{5\baselineskip}}{}{}
\pretocmd{\subsubsection}{\Needspace{4\baselineskip}}{}{}
% RULE-006: tolerancia tipografica.
\tolerance=2000 \emergencystretch=3em \hbadness=4000 \vbadness=4000
% RULE-011: todo TikZ centrado y encogido si excede el ancho util.
% ERR-0168: el salto de linea entre la figura y su pie era un punto de
%   corte valido, asi que la figura podia quedar al pie de una pagina y
%   su pie al principio de la siguiente --visto en md/lesson_02--. Una
%   figura sin pie no se entiende y un pie sin figura no dice nada, asi
%   que van juntos o van los dos a la pagina siguiente. `\Needspace*`
%   pide el hueco y `\nopagebreak[4]` prohibe el corte de dentro; no se
%   usa un `minipage` para que una figura muy alta no desborde en
%   silencio.
\newcommand{\nbtikz}[2]{%
  \par\medskip\Needspace*{6\baselineskip}\begin{center}
    \resizebox{\ifdim\width>\linewidth \linewidth\else \width\fi}{!}{#1}\\[4pt]
    \nopagebreak[4]{\small\itshape #2}
  \end{center}\par\medskip}

% =====================================================================
%  COLORES CORPORATIVOS NEBULA
% =====================================================================
\definecolor{NebulaDark}{HTML}{050B14}
\definecolor{NebulaNavy}{HTML}{0B1F3A}
\definecolor{NebulaNavyTwo}{HTML}{102A4C}
\definecolor{NebulaGold}{HTML}{C9A646}
\definecolor{NebulaCyan}{HTML}{45D6E8}
\definecolor{NebulaWhite}{HTML}{F5F7FA}
\definecolor{NebulaMuted}{HTML}{B8C4D6}
\definecolor{NebulaLight}{HTML}{F6F8FB}
\definecolor{NebulaSoft}{HTML}{EEF3F8}

\definecolor{NebulaDefinition}{HTML}{0B1F3A}
\definecolor{NebulaDefinitionBack}{HTML}{EEF3F8}
\definecolor{NebulaIdea}{HTML}{00AFCB}
\definecolor{NebulaIdeaBack}{HTML}{EAFBFD}
\definecolor{NebulaNote}{HTML}{6B7A90}
\definecolor{NebulaNoteBack}{HTML}{F2F5F8}
\definecolor{NebulaExample}{HTML}{2563A8}
\definecolor{NebulaExampleBack}{HTML}{EEF6FF}
\definecolor{NebulaStrategy}{HTML}{5B5FC7}
\definecolor{NebulaStrategyBack}{HTML}{F1F1FF}
\definecolor{NebulaWarning}{HTML}{C9A646}
\definecolor{NebulaWarningBack}{HTML}{FFF8E3}
\definecolor{NebulaError}{HTML}{B23A48}
\definecolor{NebulaErrorBack}{HTML}{FFF0F2}
\definecolor{NebulaCheck}{HTML}{2E7D5B}
\definecolor{NebulaCheckBack}{HTML}{EEF8F1}

\newcommand{\NebulaCoverSectionStyle}{\sffamily\bfseries\large}
\newcommand{\NebulaCoverBodyStyle}{\sffamily\small}
\newcommand{\NebulaCoverBoxStyle}{\sffamily\fontsize{8.2}{9.8}\selectfont}
\newcommand{\NebulaCoverProgramStyle}{\sffamily\fontsize{8.2}{9.8}\selectfont}
\newcommand{\NebulaCoverProgramTitleStyle}{\sffamily\bfseries\fontsize{8.2}{9.8}\selectfont}

% =====================================================================
%  FOOTER OBLIGATORIO
% =====================================================================
\pagestyle{fancy}
\fancyhf{}
% ---------------------------------------------------------------------
% ERR-0125 - LOS TEXTOS FIJOS DE LA PORTADA ESTABAN SOLO EN CASTELLANO.
%   Las 69 lecciones inglesas llevaban impresa la portada en espanol
%   ("Tutoria personalizada", "Acompanamiento adaptado al nivel...").
%   Ningun control lo veia: no esta en ningun JSON de leccion, esta en
%   la plantilla, y los controles de idioma miran el contenido. Ahora
%   cada cadena fija tiene sus dos versiones y manda \nebulalang.
% ---------------------------------------------------------------------
\newcommand{\nbL}[2]{\ifdefstring{\nebulalang}{en}{#2}{#1}}
\fancyfoot[C]{\footnotesize \nbL{Propiedad de Nebula.}{Property of Nebula.} Instagram: nebula.lambda. Web: https://metodonebula.es/}
\fancyfoot[R]{\footnotesize \thepage}
\renewcommand{\headrulewidth}{0pt}
\renewcommand{\footrulewidth}{0.3pt}

\titleformat{\section}{\sffamily\Large\bfseries\color{NebulaNavy}}{\thesection}{0.7em}{}
\titleformat{\subsection}{\sffamily\large\bfseries\color{NebulaNavyTwo}}{\thesubsection}{0.7em}{}
\titleformat{\subsubsection}{\sffamily\normalsize\bfseries\color{NebulaNavyTwo}}{\thesubsubsection}{0.7em}{}
\titlespacing*{\section}{0pt}{1.5em}{0.7em}

% =====================================================================
%  SISTEMA FIJO DE CAJAS SEMANTICAS
% =====================================================================
\tcbset{
  nebulaBase/.style={
    enhanced, breakable, arc=2mm, boxrule=0.65pt,
    left=9pt, right=9pt, top=9pt, bottom=9pt,
    before skip=10pt, after skip=10pt,
    fonttitle=\sffamily\bfseries, coltitle=NebulaNavy,
    attach boxed title to top left={xshift=7pt,yshift=-2mm},
    boxed title style={enhanced,arc=1.5mm,boxrule=0pt,left=5pt,right=5pt,top=2pt,bottom=2pt}
  }
}
\newtcolorbox{definitionBox}[1][]{nebulaBase,colback=NebulaDefinitionBack,colframe=NebulaDefinition,
  colbacktitle=NebulaDefinitionBack,title=\textbf{\nbIfES{Definición}{Definition}},#1}
\newtcolorbox{ideaBox}[1][]{nebulaBase,colback=NebulaIdeaBack,colframe=NebulaIdea,
  colbacktitle=NebulaIdeaBack,title=\textbf{\nbIfES{Idea clave}{Key idea}},#1}
\newtcolorbox{noteBox}[1][]{nebulaBase,colback=NebulaNoteBack,colframe=NebulaNote,
  colbacktitle=NebulaNoteBack,title=\textbf{\nbIfES{Nota}{Note}},#1}
\newtcolorbox{exampleBox}[1][]{nebulaBase,colback=NebulaExampleBack,colframe=NebulaExample,
  colbacktitle=NebulaExampleBack,title=\textbf{\nbIfES{Ejemplo guiado}{Worked example}},#1}
\newtcolorbox{strategyBox}[1][]{nebulaBase,colback=NebulaStrategyBack,colframe=NebulaStrategy,
  colbacktitle=NebulaStrategyBack,title=\textbf{\nbIfES{Estrategia}{Strategy}},#1}
\newtcolorbox{warningBox}[1][]{nebulaBase,colback=NebulaWarningBack,colframe=NebulaWarning,
  colbacktitle=NebulaWarningBack,title=\textbf{\nbIfES{Cuidado}{Take care}},#1}
\newtcolorbox{errorBox}[1][]{nebulaBase,colback=NebulaErrorBack,colframe=NebulaError,
  colbacktitle=NebulaErrorBack,title=\textbf{\nbIfES{Error frecuente}{Common mistake}},#1}
\newtcolorbox{checkBox}[1][]{nebulaBase,colback=NebulaCheckBack,colframe=NebulaCheck,
  colbacktitle=NebulaCheckBack,title=\textbf{\nbIfES{Comprobación}{Check}},#1}

% Subtitulo interno dentro de una caja (sustituye al titulo variable).
\newcommand{\nbboxhead}[1]{{\sffamily\bfseries\color{NebulaNavyTwo}#1}\par\smallskip}
\newcommand{\nbwhy}[1]{{\small\itshape\color{NebulaNote}#1}}
\newcommand{\solution}{\par\smallskip\textbf{\nbIfES{Solución}{Solution}.}\ }

% Listas de ejercicios
\newlist{ex}{enumerate}{2}
\setlist[ex,1]{label=\textbf{\arabic*.},leftmargin=1.9em,itemsep=0.55em}
\setlist[ex,2]{label=\textbf{(\alph*)},leftmargin=2em,itemsep=0.15em}
\newlist{ans}{enumerate}{2}
\setlist[ans,1]{label=\textbf{\arabic*.},leftmargin=1.9em,itemsep=0.3em}
\setlist[ans,2]{label=\textbf{(\alph*)},leftmargin=2em,itemsep=0.05em}
\newlist{steps}{enumerate}{1}
% RULE-039: itemsep/topsep subidos (antes 0.2em/0.2em). El razonamiento de
% un ejemplo resuelto es la parte que más queja de "carga visual" recibía:
% los pasos iban pegados unos a otros como un bloque único. Más espacio
% entre pasos = más páginas, pero cada paso se lee como una idea aparte.
\setlist[steps]{label=\textbf{\arabic*.},leftmargin=1.7em,itemsep=0.55em,topsep=0.4em}
\newcommand{\exlevel}[1]{\subsection*{#1}}

% RULE-023: las tarjetas van en una lista PARTIBLE, nunca en una tabla.
\newenvironment{flashlist}
  {\begin{list}{}{\setlength{\leftmargin}{0pt}\setlength{\itemindent}{0pt}%
     \setlength{\labelwidth}{0pt}\setlength{\labelsep}{0pt}%
     \setlength{\itemsep}{0.65em}\setlength{\parsep}{0pt}\setlength{\topsep}{0.3em}}}
  {\end{list}}
% ERR-0155: `\leftskip` sangra el texto pero NO estrecha las cajas. El
%   ancho util dentro de la tarjeta es `\linewidth-1.2em`, mientras que
%   `\linewidth` seguia valiendo lo de fuera, asi que `\nbtable` --que
%   encoge "hasta \linewidth" por RULE-001/022-- dejaba la tabla
%   exactamente 1.2em fuera del margen (13.14pt a \small, medido igual
%   en cuatro lecciones de tres unidades distintas). Se corrige
%   `\linewidth`, no la regla: asi tablas, ecuaciones y figuras dentro
%   de una tarjeta se calculan todas contra el ancho que de verdad hay.
%   `\leftskip` se mantiene --y no un minipage-- para que una respuesta
%   larga siga pudiendo partir de pagina.
\newcommand{\flashcard}[2]{\item\relax{\textbf{#1}}\par\nopagebreak
  {\leftskip=1.2em\relax\advance\linewidth by -1.2em\relax\small #2\par}}

\newcommand{\mcqopts}[1]{\par\nopagebreak\smallskip
  \begin{minipage}{\linewidth}
  \begin{enumerate}[label=\textbf{\Alph*)},leftmargin=2.2em,itemsep=0.1em,topsep=0.1em,nosep]
  #1
  \end{enumerate}\end{minipage}\par\smallskip}

% RULE-007: marcador que siempre compila cuando falta una figura.
\newcommand{\assetplaceholder}[2]{\par\medskip\begin{center}
  \fcolorbox{NebulaNote}{NebulaNoteBack}{\parbox{0.86\linewidth}{\centering\small
  \textbf{\nbIfES{Figura}{Figure}}\\[2pt]\itshape #2}}\end{center}\par\medskip}
\newcommand{\nbfigure}[3]{\par\medskip\begin{center}
  \includegraphics[width=#2\linewidth,max height=0.42\textheight,keepaspectratio]{#1}\\[3pt]
  {\small\itshape #3}\end{center}\par\medskip}

% Estilos TikZ comunes
\tikzset{
  nbaxis/.style={->,>=Stealth,gray!70,thick},
  nbline/.style={NebulaNavy,thick},
  nbdash/.style={NebulaNavy!60,dashed,thick},
  nbvec/.style={->,>=Stealth,NebulaError,very thick},
  nbvecb/.style={->,>=Stealth,NebulaCheck,very thick},
  nbmass/.style={circle,fill=NebulaNavy,inner sep=2pt},
  nbpoint/.style={circle,fill=black,inner sep=1.2pt},
  nblabel/.style={font=\small},
  nbangle/.style={draw=NebulaGold,thick,angle radius=8mm},
  nbfill/.style={fill=NebulaNavy!10},
}

% ERR-0110 / RULE-024 - ALIAS DE PALETA.
%   Al cambiar la plantilla a la paleta Nebula, dos lecciones dejaron de
%   compilar porque su TikZ usaba nombres de color de la version anterior.
%   La correccion no es editar las lecciones: es que la plantilla MANTENGA
%   los alias. Cambiar la paleta nunca debe romper contenido ya escrito.
\colorlet{warnR}{NebulaErrorBack}
\colorlet{tipG}{NebulaCheckBack}
\colorlet{exY}{NebulaWarningBack}
\colorlet{intuO}{NebulaWarningBack}
\colorlet{diagT}{NebulaIdeaBack}
\colorlet{chalPB}{NebulaStrategyBack}
\definecolor{navy}{HTML}{0B1F3A}
\definecolor{extR}{HTML}{B23A48}
\definecolor{coreG}{HTML}{2E7D5B}
\definecolor{greyL}{HTML}{F2F5F8}
\definecolor{defB}{HTML}{EEF3F8}
\definecolor{chalP}{HTML}{F1F1FF}

% =====================================================================
%  PORTADA FIJA NEBULA
%  \NebulaCoverFixed{tema}{titulo}{subtitulo}{tipo de documento}
% =====================================================================
\newcommand{\NebulaCoverFixed}[4]{%
\begin{titlepage}
\thispagestyle{empty}
\begin{tikzpicture}[remember picture,overlay]
\fill[NebulaDark] (current page.south west) rectangle (current page.north east);
\fill[NebulaNavy]
  (current page.north west) -- ($(current page.north east)+(0,-2.10cm)$) --
  (current page.north east) -- (current page.north west) -- cycle;
\fill[NebulaNavyTwo, opacity=0.34]
  ($(current page.north west)+(0,-4.15cm)$) -- ($(current page.north west)+(7.7cm,0)$) --
  ($(current page.north west)+(10.6cm,-4.85cm)$) -- ($(current page.north west)+(0,-6.35cm)$) -- cycle;
\fill[NebulaNavyTwo, opacity=0.26]
  ($(current page.south east)+(0,0)$) -- ($(current page.south west)+(7.5cm,0)$) --
  ($(current page.south east)+(-0.5cm,1.65cm)$) -- cycle;
\draw[NebulaCyan, opacity=0.18, line width=0.35pt]
  ($(current page.south west)+(2.00cm,0)$) -- ($(current page.north east)+(-0.9cm,-5.25cm)$);
\draw[NebulaGold, opacity=0.14, line width=0.30pt]
  ($(current page.south west)+(5.25cm,0)$) -- ($(current page.north east)+(-3.00cm,0)$);
\node[anchor=north west, text=NebulaWhite, font=\sffamily\bfseries\small]
  at ($(current page.north west)+(1.65cm,-1.35cm)$) {\nbL{METODO NEBULA}{NEBULA METHOD}};
\node[anchor=north east, text=NebulaMuted, font=\sffamily\small]
  at ($(current page.north east)+(-1.65cm,-1.35cm)$) {\nbL{Formación · Ciencias · Tecnología}{Education · Science · Technology}};
\draw[NebulaGold, line width=0.55pt]
  ($(current page.north west)+(1.65cm,-2.20cm)$) -- ($(current page.north west)+(7.25cm,-2.20cm)$);
\node[anchor=north west, rounded corners=4pt, draw=NebulaGold, line width=0.55pt,
  fill=NebulaDark, inner xsep=14pt, inner ysep=7pt, text=NebulaGold, font=\sffamily\bfseries\small]
  at ($(current page.north west)+(1.65cm,-4.85cm)$) {#4};
\node[anchor=north west, text=NebulaCyan, font=\sffamily\bfseries\Large]
  at ($(current page.north west)+(1.65cm,-6.45cm)$) {#1};
\node[anchor=north west, text=NebulaWhite,
  font=\sffamily\bfseries\fontsize{30}{43}\selectfont, text width=16.0cm, align=left]
  at ($(current page.north west)+(1.65cm,-7.55cm)$) {#2};
\node[anchor=north west, text=NebulaGold, font=\sffamily\normalsize, text width=15.8cm, align=left]
  at ($(current page.north west)+(1.70cm,-12.25cm)$) {#3};
\draw[NebulaCyan, line width=0.65pt]
  ($(current page.north west)+(1.70cm,-13.10cm)$) -- ($(current page.north west)+(8.70cm,-13.10cm)$);
\node[anchor=north west, text=NebulaWhite, font=\sffamily\normalsize, text width=15.6cm, align=left]
  at ($(current page.north west)+(1.70cm,-13.75cm)$)
  {\nbL{Material claro y estructurado para dominar el razonamiento del tema: conceptos esenciales, ejemplos guiados, estrategias de resolución y errores frecuentes.}{Clear, structured material for mastering the reasoning behind the topic: core concepts, worked examples, problem-solving strategies and the mistakes to watch for.}};
\node[anchor=north west, text=NebulaGold, font=\NebulaCoverSectionStyle]
  at ($(current page.north west)+(1.65cm,-15.85cm)$) {\nbL{Método Nebula}{The Nebula Method}};
\node[anchor=north west, text=NebulaWhite, font=\NebulaCoverBodyStyle, text width=15.6cm, align=left]
  at ($(current page.north west)+(1.65cm,-16.60cm)$)
  {\nbL{Acompañamiento académico y técnico para estudiar con estructura, comprender de verdad y avanzar con un plan claro.}{Academic and technical support: study with structure, understand in depth, and move forward with a clear plan.}};
\node[anchor=north west, rounded corners=5pt, draw=NebulaCyan, line width=0.45pt, fill=NebulaNavy,
  fill opacity=0.20, text opacity=1, inner xsep=6pt, inner ysep=7pt, text=NebulaWhite,
  font=\NebulaCoverBoxStyle, text width=4.00cm, align=center]
  at ($(current page.north west)+(1.65cm,-17.95cm)$) {\nbL{Entender conceptos clave.}{Understand the key ideas.}};
\node[anchor=north west, rounded corners=5pt, draw=NebulaGold, line width=0.45pt, fill=NebulaNavy,
  fill opacity=0.20, text opacity=1, inner xsep=6pt, inner ysep=7pt, text=NebulaWhite,
  font=\NebulaCoverBoxStyle, text width=4.60cm, align=center]
  at ($(current page.north west)+(6.45cm,-17.95cm)$) {\mbox{\nbL{Preparar pruebas con estrategia.}{Prepare for exams with a strategy.}}};
\node[anchor=north west, rounded corners=5pt, draw=NebulaCyan, line width=0.45pt, fill=NebulaNavy,
  fill opacity=0.20, text opacity=1, inner xsep=6pt, inner ysep=7pt, text=NebulaWhite,
  font=\NebulaCoverBoxStyle, text width=4.00cm, align=center]
  at ($(current page.north west)+(11.85cm,-17.95cm)$) {\nbL{Avanzar con ruta definida.}{Progress along a clear path.}};
\node[anchor=north west, text=NebulaGold, font=\NebulaCoverSectionStyle]
  at ($(current page.north west)+(1.65cm,-19.55cm)$) {\nbL{Programas}{Programmes}};
\node[anchor=north west, rounded corners=5pt, draw=NebulaCyan, line width=0.45pt, fill=NebulaNavy,
  fill opacity=0.22, text opacity=1, inner xsep=7pt, inner ysep=7pt, text=NebulaWhite,
  font=\NebulaCoverProgramStyle, text width=4.60cm, minimum height=1.30cm, align=left]
  at ($(current page.north west)+(1.65cm,-20.35cm)$)
  {{\NebulaCoverProgramTitleStyle \nbL{Tutoría personalizada}{One-to-one tutoring}}\\\nbL{Acompañamiento adaptado al nivel\\y ritmo del alumno.}{Support matched to each\\student's level and pace.}};
\node[anchor=north west, rounded corners=5pt, draw=NebulaGold, line width=0.45pt, fill=NebulaNavy,
  fill opacity=0.22, text opacity=1, inner xsep=7pt, inner ysep=7pt, text=NebulaWhite,
  font=\NebulaCoverProgramStyle, text width=4.60cm, minimum height=1.30cm, align=left]
  at ($(current page.north west)+(7.05cm,-20.35cm)$)
  {{\NebulaCoverProgramTitleStyle \nbL{Preparación por objetivos}{Goal-driven preparation}}\\\nbL{Plan enfocado en exámenes,\\entregas y pruebas.}{A plan built around exams,\\deadlines and assessments.}};
\node[anchor=north west, rounded corners=5pt, draw=NebulaCyan, line width=0.45pt, fill=NebulaNavy,
  fill opacity=0.22, text opacity=1, inner xsep=7pt, inner ysep=7pt, text=NebulaWhite,
  font=\NebulaCoverProgramStyle, text width=4.60cm, minimum height=1.30cm, align=left]
  at ($(current page.north west)+(12.45cm,-20.35cm)$)
  {{\NebulaCoverProgramTitleStyle \nbL{Mejora y transición IT}{IT upskilling and transition}}\\\nbL{Ruta de progreso en programación,\\datos e IA.}{A route through coding,\\data and AI.}};
\node[anchor=south west, text=NebulaMuted, font=\sffamily\small, text width=17.7cm, align=left]
  at ($(current page.south west)+(1.65cm,2.05cm)$)
  {\nbL{Más información en}{More at} https://metodonebula.es/ \textcolor{NebulaGold}{|} Instagram: nebula.lambda};
\end{tikzpicture}
\vspace*{\fill}
\end{titlepage}}

%% B-02 / RULE-028 - biblioteca de figuras parametrizadas
\providecommand{\nbtpl}{}
\input{\nbtpl nebula_figures}
\input{\nbtpl nebula_waves}
\input{\nbtpl nebula_optics}
